\origpage{254}{ROBERTSON AND SEYMOUR}
\textit{Proof.} For each $i \ge 1$, let $G_i$ be as in (iii) above, and for each $e \in E_i$ let $\Delta_i(e)$ be the set of all pairings in $\bar{\mu}_i(e)$ no two edges of which cross (where $\{a, c\}$ crosses $\{b, d\}$ if $a, b, c, d \in \bar{e} - e$ are all distinct and occur in that order, or the reverse, in the natural cyclic order of $\bar{e} - e$). Then $P_i = (G_i, \mu_i, \Delta_i, \phi_i)$ is an $\Omega$-patchwork, and it is robust. By (9.1) there exist $i' > i \ge 1$ and a realizable expansion $\eta$ of $P_i$ in $P_{i'}$. Let $H$ be a realization of $P_{i'} \setminus \eta(E(G_i))$ realizing $\eta$, and let $F$ be a spanning forest of $H$. From our choice of the $\Delta_i(e)$'s and the fact that $H$ is a realization of $P_{i'} \setminus \eta(E(G_i))$, there is a drawing $\Xi$ of $F$ in $\Sigma$ such that each vertex of $F$ is represented by itself (we recall that $V(F) = V(H) = N_{i'} \subseteq \Sigma$) and each edge is drawn within some cell $e' \in E_{i'} - \eta(E_i)$. Define $\sigma(e) = \eta(e)$ $(e \in E_i)$ and for each $v \in N_i$, let $\sigma(v)$ be the component of $\Xi$ with vertex set $\eta(v)$. Then (a), ..., (g) are satisfied, as required. $\square$

\sectionhead{\textsc{References}}

1. D. Archdeacon, ``A Kuratowski Theorem for the Projective Plane,'' Thesis, Ohio State University, 1980.

2. H. Glover, P. Huneke, and C. S. Wang, 103 graphs that are irreducible for the projective plane, \textit{J. Combin. Theory Ser. B} \textbf{27} (1979), 332-370.

3. G. Higman, Ordering by divisibility in abstract algebras, \textit{Proc. London Math. Soc.} \textbf{2} (1952), 326-336.

4. T. A. Jenkyns and C. St. J. A. Nash-Williams, Counterexamples in the theory of well-quasi-ordered sets, in ``Proof Techniques in Graph Theory,'' Proceedings of the Second Ann Arbor Graph Theory Conference, February, 1968, pp. 87-91, Academic Press, New York/London 1969.

5. J. Kruskal, Well-quasi-ordering, the tree theorem, and V\'azsonyi's conjecture, \textit{Trans. Amer. Math. Soc.} \textbf{95} (1960), 210-225.

6. W. Mader, Wohlquasiordnete Klassen endlicher Graphen, \textit{J. Combin. Theory Ser. B} \textbf{12} (1972), 105-122.

7. C. St. J. A. Nash-Williams, ``On well-quasi-ordering finite trees,'' \textit{Proc. Camb. Phil. Soc.} \textbf{59} (1963), 833-835.

8. N. Robertson and P. D. Seymour, Graph minors. II. Algorithmic aspects of tree-width, \textit{J. Algorithms} \textbf{7} (1986), 309-322.

9. N. Robertson and P. D. Seymour, Graph minors. III, Planar tree-width, \textit{J. Combin. Theory Ser. B} \textbf{36} (1984), 49-64.

10. N. Robertson and P. D. Seymour, Graph minors. V. Excluding a planar graph, \textit{J. Combin. Theory Ser. B} \textbf{41} (1986), 92-114.

11. N. Robertson and P. D. Seymour, Graph minors. VIII. A Kuratowski theorem for general surfaces, \textit{J. Combin. Theory Ser. B} \textbf{48} (1990), 255-288.

12. S. G. Simpson, Non-provability of certain combinatorial properties of finite trees, in ``Harvey Friedman's Research of the Foundations of Mathematic'' (L. Harrington, M. Morley, A. Schredov, S. Simpson, Eds.), pp. 87-117, North-Holland, Amsterdam, 1985.

13. R. Thomas, A Menger-like property of tree-width. The finite case, \textit{J. Combin. Theory Ser. B}, to appear.
